\scanpage{77}\mspage{76}

\deleted{\uline{\unsure{Prop.|Déf.}}} \inserted{\uline{Prop.}} Soit

\diagram{\begin{tikzcd}[ampersand replacement=\&,row sep=0.15em,column sep=0.3em]
 G \& \ll \& F \\
 \deleted{\rotatebox{90}{$\leq$}} \& {} \& \rotatebox{90}{$\leq$} \\
 \deleted{\unsure{G'|G?}} \& {} \& F'\,\unsure{.}
\end{tikzcd}}

alors il existe \unsure{un|une} plus grande sous figure $F'$ de $G$, qui raffine $F'$ (et qu'on appelle raffinement de $F'$ induit \deleted{par le} sur la sous figure $F'$ par le raffinement $G$ de $F$, ou la sous-figure du \deleted{$G$ \unsure{induit}} \unsure{raff\textsuperscript{t}|raff} $G$ de $F$, induite par la sous-figure $F'$ de $F$). \deleted{\unsure{Elle est}} \unsure{Ses|Les} strates \deleted{(c'est $\unsure{\pm|=}$ tautologique)} sont les strates de $G$ qui raffinent $F'$. \deleted{De plus, $G$ et $F$}

\diagram{\begin{tikzcd}[ampersand replacement=\&,row sep=0.12em,column sep=0.3em]
 G \& \ll \& F \\
 \rotatebox{90}{$\leq$} \& {} \& \rotatebox{90}{$\leq$} \\
 G' \& \ll \& F'
\end{tikzcd}}

De plus, $G'$ est un $\operatorname{Inf}$ de $G$ et $F'$ pour la relation $\underset{\unsure{\cdot}}{\ll}$, en d'autres termes, si \deleted{\unsure{$X$} $G$ est} $K$ est une figure telle que $K\ll G$ et $K\ll F'$, alors $K\ll G'$.

\margin{Dire (ou poser) que $G'=\varnothing\Longleftrightarrow G$ disjoint de $F'$ (seul $\Rightarrow$ non évident\unsure{....|...})}

\deleted{\uline{\unsure{Dém|Démon}} Soit $H$ une sous-figure \deleted{de $G$ telle que} \inserted{de $H$} $H\ll F'$, alors pour tte strate $X\leq H$, on aura $X\ll F'$, donc \unsure{Si} on désigne par $G'$ la sous-figure de $G$ engendrée par des strates $Y\leq G$ telles que $Y\ll H$, alors \deleted{$\unsure{X}$} $H\leq G$ puisque}
\[
 \deleted{\displaystyle H=\mathop{\mathrm{Sup}}\limits_{\substack{X\in\deleted{\unsure{H|G}}\,\unsure{M|\mathcal M}\\X\leq H}}X\deleted{\unsure{<}}\leq\mathop{\mathrm{Sup}}\limits_{\substack{X\in M\\X\in G'}}X=G'\unsure{.|,}}
\]
\deleted{D'autre part, $G'\ll$}
